\documentclass[12pt]{article}
\usepackage[utf8]{inputenc}
\usepackage[T1]{fontenc}
\usepackage{hyperref}
\title{Aramba Language Resources for K-12 Curriculum Development}
\author{Compiled by Claude Sonnet 5.5 (claude-sonnet-5-5) for Andr\'as Kornai}
\date{}
\begin{document}
\maketitle

\section{Introduction}
Aramba (also written Arammba; alternative names include Serki, Upper Morehead and others) is a Papuan language of the Tonda branch, spoken in a few villages between the Fly River and Morehead in Western Province, Papua New Guinea. Estimates range from about 750 to 2,000 speakers (Wikipedia cites 970). The ELDP documentation project notes that the number of speakers is declining.

\section{Primary Educational Resources}
No primers or school textbooks were found in search results. The reference description says Aramba is used as a language of instruction, but I found no evidence of materials for it.

\section{Linguistic and Reference Materials}
\begin{itemize}
\item Aramba language, Wikipedia. \url{https://en.wikipedia.org/wiki/Aramba_language}
\item Arammba, Omniglot. \url{https://www.omniglot.com/writing/arammba.htm}
\item Aramba Webonary (online dictionary, introduction page). \url{https://www.webonary.org/aramba/overview/introduction/}
\item A linguistic description of Arammba, an endangered language of southern New Guinea (Cultures in Crisis project page). \url{https://cultureincrisis.org/projects/a-linguistic-description-of-arammba-an-endangered-language-of-southern-new-guinea}
\item Glottolog entry aram1253. \url{https://glottolog.org/resource/languoid/id/aram1253}
\item Arammba numerals page (MPI-EVA). \url{https://lingweb.eva.mpg.de/channumerals/Arammba.htm}
\end{itemize}

\section{Religious and Community Resources}
\begin{itemize}
\item Joshua Project language page for stk. \url{https://joshuaproject.net/index.php/languages/stk}
\end{itemize}

\section{Digital and Audio Resources}
\begin{itemize}
\item ELDP archive collection for the Arammba documentation (E. Mittag). \url{https://eldp.access.preservica.com/dk0366/}
\item PARADISEC catalogue entry holding recordings. \url{https://catalog.paradisec.org.au/repository/LSNG06/IND02}
\end{itemize}

\section{Curriculum Development Context}
The ELDP and PARADISEC recordings and the Webonary dictionary are the most substantial resources found and could supply vocabulary and texts for classroom material. I did not open the Webonary pages, so the dictionary's size and scope are unconfirmed. The small speaker community means any curriculum would need to be developed with the villages concerned.

\end{document}